\chapter{Validation Campaign}

The Reference Node has been exercised through functional tests and long-duration acquisition runs.

\section{Initial Long-Duration Test}

A representative long-duration test reached an uptime of 68,395 seconds, with 49,425 successful measurements, zero measurement failures, zero TWI errors and zero EEPROM errors.

UART receive robustness was subsequently improved through the finite-state-machine command parser used in NodeOS v0.5.3-alpha1.

\section{Extended Long-Duration Test}

A later extended acquisition run provided a substantially longer operational sample. As documented in Figure~\ref{fig:monitor_long_run}, the Reference Single-Sensor Node reached an uptime of 273792 seconds (approximately 76 hours). At the captured point the diagnostic counters reported 197589 successful measurements, zero measurement failures, zero TWI errors, zero UART receive overflows and zero EEPROM errors.

The ND280 Monitor simultaneously reported zero missing packets for the displayed acquisition sequence. This extended run provides additional evidence that the acquisition, TWI communication, EEPROM handling, command transport and monitoring chain can operate continuously over multi-day periods without the error modes observed during the earlier UART parser development phase.

These results validate the Reference Node as an engineering baseline. Final noise, offset, thermal behaviour, absolute calibration and distributed-bus performance remain to be qualified on the Eight-Sensor Multiplexing Node and on the final detector hardware.


\section{Preliminary Stability Characterization}
\label{sec:preliminary_stability}

Before a controlled metrological calibration campaign is performed, data acquired with the Reference Single-Sensor Node already provide a useful first characterization of the observed short-term stability of the complete acquisition chain. The results presented in this section are based on a continuous dataset containing 34,418 magnetic-field measurements. They should be interpreted as \emph{preliminary system-level observations}: they include the sensor, embedded averaging and readout chain, local environmental conditions and any real variation of the ambient magnetic field. They are therefore not measurements of the intrinsic TMAG5273A2 resolution or of its absolute accuracy.

\subsection{Observed field distributions}

Figure~\ref{fig:prelim_bxyz_distributions} shows the distributions of the three measured field components. After rebinning, Gaussian functions were fitted as a compact first description of the observed widths. The fitted standard deviations are approximately \SI{0.0427}{mT}, \SI{0.0464}{mT} and \SI{0.0314}{mT} for $B_x$, $B_y$ and $B_z$, respectively. The corresponding fitted means are approximately \SI{-0.0944}{mT}, \SI{-0.2921}{mT} and \SI{0.1228}{mT}.

\begin{figure}[htbp]
  \centering
  \includegraphics[width=0.98\textwidth]{VV01_Bxyz_Distributions_Gaussian_Fits.pdf}
  \caption{Preliminary distributions of the three magnetic-field components measured with the Reference Single-Sensor Node over 34,418 samples. Gaussian fits are superimposed after rebinning and are used only as a first characterization of the observed distribution widths.}
  \label{fig:prelim_bxyz_distributions}
\end{figure}

Table~\ref{tab:preliminary_field_stability} summarizes the measured distribution means and RMS values together with the parameters obtained from the rebinned Gaussian fits.

\begin{table}[htbp]
\centering
\caption{Preliminary statistical characterization of the three magnetic-field components measured during the Reference Node stability run. Gaussian parameters are reported as descriptive quantities only and shall not be interpreted as measurements of the intrinsic TMAG5273A2 resolution.}
\label{tab:preliminary_field_stability}
\small
\begin{tabular}{lcccc}
\hline
Component &
\begin{tabular}[c]{@{}c@{}}Distribution\\mean [mT]\end{tabular} &
\begin{tabular}[c]{@{}c@{}}Distribution\\RMS [mT]\end{tabular} &
\begin{tabular}[c]{@{}c@{}}Gaussian\\mean [mT]\end{tabular} &
\begin{tabular}[c]{@{}c@{}}Gaussian\\$\sigma$ [mT]\end{tabular} \\
\hline
$B_x$ & $-0.09286$ & $0.04438$ & $-0.09438 \pm 0.00024$ & $0.04272 \pm 0.00020$ \\
$B_y$ & $-0.2894$  & $0.04541$ & $-0.2921 \pm 0.0003$  & $0.04643 \pm 0.00017$ \\
$B_z$ & $+0.1223$  & $0.03337$ & $+0.1228 \pm 0.0002$  & $0.03141 \pm 0.00015$ \\
\hline
\end{tabular}
\end{table}

The fitted Gaussian widths correspond to approximately \SI{43}{\micro\tesla}, \SI{46}{\micro\tesla} and \SI{31}{\micro\tesla} for the $B_x$, $B_y$ and $B_z$ components, respectively. These values characterize the observed short-term spread of the complete acquisition chain rather than the intrinsic resolution of the magnetic sensor.

The Gaussian description is not statistically adequate for the complete distributions: the fitted $\chi^2/\mathrm{ndf}$ values are approximately 3968/17, 3570/15 and 2519/16 for $B_x$, $B_y$ and $B_z$, respectively. Consequently, the fitted $\sigma$ values must not be interpreted as intrinsic sensor resolutions. The non-Gaussian structure may contain contributions from quantization, slow drift, temperature dependence, environmental magnetic-field variations and other effects of the complete measurement chain. This observation directly motivates the controlled characterization programme described later in this chapter.

\subsection{Temporal behaviour}

The same dataset was inspected as a function of acquisition time. Figure~\ref{fig:prelim_bxyz_time} overlays the three measured field components throughout the run. No correction for temperature or slow environmental variation has been applied in this representation.

\begin{figure}[htbp]
  \centering
  \includegraphics[width=0.96\textwidth]{VV02_B_Components_vs_Time.pdf}
  \caption{Preliminary time evolution of $B_x$, $B_y$ and $B_z$ during the stability dataset. The plot represents the measured values before a dedicated temperature-dependent or metrological correction is applied.}
  \label{fig:prelim_bxyz_time}
\end{figure}

This view is complementary to the one-dimensional distributions because it allows abrupt excursions, slow drifts and time-correlated structures to be distinguished from purely point-to-point fluctuations. A quantitative stability specification will only be assigned after equivalent measurements have been repeated under controlled magnetic-field and temperature conditions.

\subsection{Preliminary temperature dependence}

The recorded TMAG5273A2 temperature allows a first investigation of whether changes in the reported magnetic-field components are correlated with sensor temperature. Figure~\ref{fig:prelim_bxyz_temperature} shows profiles of the mean $B_x$, $B_y$ and $B_z$ values versus temperature for the same dataset. The sampled temperature interval is approximately \SIrange{26}{32}{\celsius}, with a mean close to \SI{28}{\celsius}.

\begin{figure}[htbp]
  \centering
  \includegraphics[width=0.96\textwidth]{VV03_B_Profiles_vs_Temperature.pdf}
  \caption{Profile of the mean measured magnetic-field components versus TMAG5273A2 temperature. The profiles are intended to reveal possible temperature-correlated behaviour in the complete acquisition chain; they do not by themselves determine the intrinsic temperature coefficients of the sensor.}
  \label{fig:prelim_bxyz_temperature}
\end{figure}

Any apparent slope in these profiles must be interpreted cautiously. In an uncontrolled run, temperature, elapsed time and the ambient magnetic field are not independent variables. A correlation between $B$ and temperature can therefore arise from sensor response, from a real environmental field change correlated with time, or from both. The purpose of this preliminary analysis is to identify effects that must subsequently be separated in a controlled calibration campaign, not to derive final temperature-correction coefficients.

Taken together, the distributions, time evolution and temperature profiles establish a useful pre-calibration baseline for the Reference Node. They also demonstrate why long-duration stability and metrological characterization are complementary: operational continuity alone does not establish measurement accuracy, while a calibration model must be supported by evidence that it remains stable during realistic operation.

\section{Calibration and Metrological Characterization Strategy}

\subsection{Scope of the calibration programme}

The ND280 Magnetic Field Monitoring System is intended primarily as a distributed monitoring instrument rather than as an absolute magnetic-field mapping system. Its calibration strategy must therefore be proportional to this objective. The system does not need to reproduce the metrological performance of dedicated high-precision Hall mapping probes; however, calibration cannot be neglected because sensor-to-sensor differences, offsets, temperature dependence and cross-axis effects could otherwise generate apparent field variations that are unrelated to the detector magnet.

A useful reference is the calibration and characterization procedure developed for a dedicated CERN 3D Hall sensor card \cite{MessnerHallCalibration}. That work identifies offset, temperature dependence, non-linearity, planar Hall effects and mechanical misalignment as relevant contributors to the final measurement error, and demonstrates the value of a controlled calibration procedure for compensating such effects. The authors report mean residuals of approximately $9.15\times10^{-5}$~T in field magnitude and 0.97~mrad in angle over a 0.1--2~T range across ten calibrated units. These values define the performance of a substantially more demanding mapping application and are not adopted as requirements for the ND280 monitor; they nevertheless provide a useful example of how calibration accuracy, temperature stability and production reproducibility can be quantified.

For the ND280 system, the principal metrological objectives are instead:

\begin{itemize}[leftmargin=*]
\item repeatable measurement of temporal changes in the magnetic field;
\item comparability of measurements from different TMAG5273A2 sensors and acquisition boards;
\item controlled dependence of the reported field on sensor temperature;
\item sufficient absolute accuracy to identify incorrect installation, configuration or gross field deviations;
\item traceable calibration parameters that remain associated with the physical sensor throughout operation and maintenance.
\end{itemize}

\subsection{Proposed calibration model}

The present NodeOS implementation supports per-axis offset and gain parameters stored in non-volatile memory. This representation is intentionally simple and remains appropriate for early bring-up and engineering tests. The calibration infrastructure will, however, be designed so that it can evolve without changing the higher-level acquisition protocol.

The natural next step is a vector calibration model of the form
\begin{equation}
\vec{B}_{\mathrm{cal}} = \mathbf{M}\left(\vec{B}_{\mathrm{raw}}-\vec{O}\right),
\label{eq:nd280_vector_calibration}
\end{equation}
where $\vec{O}$ is a three-component offset vector and $\mathbf{M}$ is a $3\times3$ correction matrix. The diagonal terms of $\mathbf{M}$ describe axis sensitivities, while the off-diagonal terms can absorb first-order cross-axis coupling and small orientation errors. A temperature-dependent extension may be introduced after characterization, for example by allowing the offset vector and selected matrix coefficients to vary around a reference temperature $T_0$.

The complexity of the final model will be determined experimentally. The CERN Hall-card study uses a considerably more sophisticated response model based on spherical harmonics because it addresses discrete Hall elements, mechanical alignment and a very wide field range \cite{MessnerHallCalibration}. Such a model would be unnecessarily complex for the present system unless the TMAG5273A2 characterization demonstrates residual non-linearities that cannot be corrected with a simpler matrix-based approach.

\subsection{Field-range characterization}

The dedicated Hall-card study shows that calibration parameters may vary with magnetic-field strength and describes a multi-field procedure in which calibration coefficients determined at different field magnitudes are interpolated across the operating range \cite{MessnerHallCalibration}. This observation is relevant even though the ND280 monitoring system operates in a much narrower field region centered near \SI{0.2}{T}.

The first Rev.A characterization campaign should therefore include several reference field values around the expected operating region rather than relying on a single-point calibration. A preliminary set of characterization points may include zero field and representative values around \SI{0.1}{T}, \SI{0.2}{T} and \SI{0.25}{T}, subject to the capabilities of the available reference magnet and instrumentation. These values are not yet calibration requirements; they define an initial measurement programme intended to determine whether a single linear correction is sufficient in the ND280 operating range.

If the residuals show a significant field-dependent structure, the calibration model can be extended either by piecewise interpolation of coefficients or by an alternative non-linear correction. If no significant improvement is observed, the simpler model will be retained.

\subsection{Temperature characterization}

Temperature must be treated as a measured calibration variable rather than only as an environmental diagnostic. This requirement is facilitated by the TMAG5273A2, which already provides local temperature information for every sensor.

The CERN calibration study demonstrates why this is important: the residual error in reconstructed field magnitude changes measurably between the calibration temperatures and the temperature extremes, even after application of a temperature correction model \cite{MessnerHallCalibration}. The ND280 characterization campaign should therefore repeat representative field measurements at several temperatures spanning the expected detector operating range. The objective is not to reproduce the temperature-controlled apparatus of the reference study, but to quantify the temperature coefficient and determine whether explicit compensation is required for the desired monitoring precision.

\subsection{Sensor-to-sensor and board-to-board reproducibility}

A distributed system containing potentially hundreds of sensors cannot rely on the behaviour of a single characterized device. The calibration campaign must therefore quantify the spread between TMAG5273A2 devices and between independently assembled acquisition boards.

The reference Hall-card work calibrated and tested ten separate cards and used the dispersion of the reconstructed field residuals as a measure of production and calibration reproducibility \cite{MessnerHallCalibration}. A similar principle should be adopted for the ND280 system. After Rev.A bring-up, multiple boards should be exposed to identical magnetic-field and temperature conditions and their residual distributions compared. This will determine whether one generic calibration model is sufficient, whether coefficients must be determined sensor-by-sensor, and how much of the measurement uncertainty originates from production variability.

\subsection{Calibration data management}

Calibration data shall remain permanently associated with the physical sensor. The current EEPROM infrastructure already provides the basis for persistent storage, but the data format should evolve from the present simple gain-and-offset representation to a versioned calibration record. A future calibration record should contain, as appropriate:

\begin{itemize}[leftmargin=*]
\item calibration-model version;
\item sensor and board identifiers;
\item three-axis offset vector;
\item three-dimensional correction matrix;
\item temperature coefficients and reference temperature;
\item calibration date and calibration-range identifier;
\item validation or quality flags;
\item integrity information such as a CRC.
\end{itemize}

Versioning is important because it allows NodeOS to retain backward compatibility with the existing Reference Node while supporting more complete calibration models on later hardware revisions.

\subsection{In-situ verification}

Laboratory calibration and detector operation serve different purposes. Laboratory calibration determines the relationship between sensor output and a known magnetic field under controlled conditions. Once installed in ND280, the monitoring network should instead perform periodic consistency checks intended to reveal drift, installation problems or anomalous sensors.

In-situ checks may include comparisons between nearby sensors, repeated measurements at reproducible magnet operating points, verification of zero-field offsets when such operating conditions are available, and comparison with the established ND280 magnetic-field map. These procedures do not replace laboratory calibration; they provide long-term verification that the calibration remains applicable throughout detector operation.

\subsection{Calibration work package for Rev.A}

The calibration programme will therefore be introduced as a dedicated work package following the electrical bring-up of the Eight-Sensor Multiplexing Node. The initial objective is characterization rather than immediate implementation of a complex correction algorithm. Raw data from multiple sensors, field values and temperatures will first be collected and analysed offline. Only after the dominant error terms have been identified will the corresponding correction model be frozen and incorporated into NodeOS and the EEPROM calibration format.

This staged approach is consistent with the overall engineering philosophy of the project: measurement evidence determines the complexity of the implementation, and each additional calibration feature must demonstrate a measurable improvement before becoming part of the reference architecture.
